\uuid{BObq}
\titre{Application linéaire}
\auteur{hanani ; mzari}
\organisation{exo7}
\licence{CC BY-NC-SA 4.0}
\datecreate{2020-01-20}
\niveau{L1}
\module{Algèbre}
\chapitre{Application linéaire}
\sousChapitre{Définition}
\difficulte{3}
\theme{Applications linéaires, Application linéaire}
\exo7qcmid{lille-2.applications-lineaires.009}

\contenu{
\begin{question}
\qclassification{107.01}
\qtags{motcle=Application linéaire}
On considère les applications suivantes :
$$\begin{array}{rccc}f:&\Rr^3&\to&\Rr^2\\
& (x,y,z)&\to &(z,x+ax^2) \,  \end{array}  \quad \mbox{et} \quad  \begin{array}{rccc}g:&\Rr^3&\to&\Rr^3\\
& (x,y,z)&\to &(z+a\sin x, y+be^x, c|x|+1).\end{array} $$
où $a,b$ et $c$ sont des réels. Quelles sont les assertions vraies ?
\begin{answers}
\bad{Pour tout $a\in \Rr$, $f$ est une application linéaire}.
\good{$f$ est une application linéaire si et seulement si $a=0$}.
\bad{$g$ est une application linéaire si et seulement si $a=b=c=0$}.
\good{Pour tous $a,b, c \in \Rr$, $g$ n'est pas une application linéaire}.
\end{answers}
\begin{explanations} On vérifie que $f$ est linéaire si et seulement si $a=0$ et que $g$ n'est pas linéaire pour
tous réels $a,b$ et $c$.
\end{explanations}
\end{question}
}
